% Conclusions for the FedXR-GAN project
\section{Conclusions}
\label{sec:conclusions}

This project shows that a practical pipeline for pneumonia
classification can be built from two complementary components. First a
DCGAN \cite{Radford2015DCGAN} was trained to augment the minority
\texttt{NORMAL} class. The generated set reduced class imbalance and
achieved a measured FID around 25 \cite{Heusel2017FID} which indicates
reasonable visual fidelity for downstream use.

Second a compact CNN classifier trained in a federated setup
\cite{McMahan2017FedAvg} reached high discrimination on the evaluation
set. Over 50 communication rounds the evaluation AUC stabilized near
0.99 and the evaluation loss reached values near 0.11. These results
demonstrate that federated training can match centralized performance
while keeping data local to simulated clients in IID data split
scenarios. 

Taken together the GAN augmentation and the federated CNN form a practical approach for sensitive medical imaging tasks. The pipeline improves data availability and preserves privacy while delivering strong classification metrics. At the same time there remain clear directions for improvement. The FID can be lowered by longer GAN training and richer data augmentation. The classifier would benefit from additional cross site validation and adversarial robustness checks. Finally real world deployment requires formal privacy auditing and fairness analysis across demographic groups.

Overall the work validates the feasibility of combining synthetic data generation with federated learning for medical image classification. Figures~\ref{fig:fed_curves} and~\ref{fig:client_curves} show the global and per client training behaviour reported in the experiments. The experiments provide a reproducible baseline and a set of artifacts to continue development. Numeric logs are available in `results/federated\_metrics.csv` and `fid\_score\_results.txt` in this repository. The full code and generated artifacts are published on the project GitHub repository \url{https://github.com/cyberpabs/FedXR-GAN}.
